\documentclass{beamer}
\usetheme{default}

% -----------------------------------------------------------------------
% AGU-specific directives (comment sentinels; invisible to Beamer)
% -----------------------------------------------------------------------
%% AGU_EMAIL: author@university.edu
%\journalname{Geophysical Research Letters}

\usepackage{natbib}
\usepackage{amsmath}
\usepackage{amssymb}
\usepackage{bm}

\title[Short Title]{Full Test Title for Conversion Scripts}
\author[]{Author One\inst{1} and Author Two\inst{2} and Author Three\inst{1}}
\institute[]{%
  \inst{1} Department of Atmospheric Science, Test University, Boulder, CO, USA
  \and
  \inst{2} Institute of Oceanography, Second University, Seattle, WA, USA
}

\begin{document}

% -----------------------------------------------------------------------
% Abstract
% -----------------------------------------------------------------------
\section{Abstract}
\begin{frame}{Abstract}
This study tests all features of the Beamer-to-manuscript conversion
pipeline. We verify metadata extraction, citation remapping, equation
wrapping, figure and table formatting, list flattening, and supplemental
information handling.
\end{frame}

% -----------------------------------------------------------------------
% Key Points  (AGU: consumed into \begin{keypoints}; AMS: passthrough section)
% -----------------------------------------------------------------------
\section*{Key points}
\begin{frame}{Key Points}
\begin{itemize}
  \item Conversion pipeline correctly extracts metadata from Beamer preamble
  \item All natbib citation commands are remapped to apacite for AGU output
  \item Supplemental figures and tables are counted and listed automatically
\end{itemize}
\end{frame}

% -----------------------------------------------------------------------
% Plain Language Summary  (AGU only: consumed into preamble PLS slot)
% -----------------------------------------------------------------------
\section*{Plain Language Summary}
\begin{frame}{Plain Language Summary}
Scientific papers often begin as slide presentations. This tool automates
the conversion of slides to journal manuscript format, reducing manual
reformatting effort.
\end{frame}

% -----------------------------------------------------------------------
% Introduction
% -----------------------------------------------------------------------
\section{Introduction}

\begin{frame}{Background}
% Citation variants: citet, citep, citealp, citeauthor, citeyear
The Madden--Julian Oscillation \citep{madden1971} is a tropical weather
pattern with a 30--90 day period. \citet{wheeler2004} developed the RMM
index. See also \citealp{zhang2005} and \citeauthor{lau2012}
(\citeyear{lau2012}).
\end{frame}

\begin{frame}{Motivation}
\begin{itemize}
  % One-arg citep with postnote:
  \item Forecast skill degrades rapidly \citep[e.g.,]{vitart2017}
  % Two-arg citet (pre + post):
  \item Neural models show promise \citet[see][their Fig.~3]{smith2020}
  % Two-arg citet with empty prenote:
  \item Further evidence \citet[][p.~5]{jones2019}
  % Two-arg citep (pre + post):
  \item See also \citep[cf.][Table~2]{brown2021}
  % Beamer formatting that should be stripped:
  \item \alert{Highlighted text} and \textcolor{red}{colored text} are removed
  \item \only<2>{Overlay-only content is kept but tag is stripped}
  \item \structure{Structural markup} is unwrapped
\end{itemize}
\end{frame}

% -----------------------------------------------------------------------
% Methods
% -----------------------------------------------------------------------
\section{Methods}

\begin{frame}{Equations}
Anomalies are computed as
\begin{equation}
  u'(\lambda, t) = u(\lambda, t) - \bar{u}(\lambda, t),
  \label{eq:anomaly}
\end{equation}
where $\bar{u}$ is the climatological mean.\pause

Display math with bracket notation:
\[
  \mathrm{PC1}(t) = \int_{-15}^{15} \bm{e}_1 \cdot \bm{u}'\,d\lambda
\]

Aligned system:
\begin{align}
  x &= r\cos\theta \label{eq:x}\\
  y &= r\sin\theta \label{eq:y}
\end{align}
\end{frame}

\begin{frame}{Beamer Formatting}
% block → bold title
\begin{block}{Block Title}
Block environments are converted to a bold title.
\end{block}
\vspace{1em}
% center → unwrap
\begin{center}
Centered content is unwrapped.
\end{center}
\medskip
% columns → strip tags, keep content
\begin{columns}
\column{0.5\textwidth}
Left column content.
\column{0.5\textwidth}
Right column content.
\end{columns}
% font size commands should be stripped
\small Small text here. \normalsize Normal again.
\end{frame}

% -----------------------------------------------------------------------
% Results
% -----------------------------------------------------------------------
\section{Results}

\begin{frame}{Main Figure}
% \fig{} macro → \includegraphics via preprocess_body
\fig{test_figure_1.pdf}
\captionof{figure}{Spatial pattern of EOF-1 for U850 (red) and U200 (green).}
\end{frame}

\begin{frame}{Two Figures Side by Side}
% Consecutive \includegraphics → grouped into one figure environment
\includegraphics[width=0.48\textwidth]{test_figure_2a.pdf}
\includegraphics[width=0.48\textwidth]{test_figure_2b.pdf}
\caption{Left: PC1 power spectrum. Right: PC2 power spectrum.}
\label{fig:spectra}
\end{frame}

\begin{frame}{Summary Table}
\begin{tabular}{lcc}
\hline
Metric         & Forecast & Observed \\
\hline
Mean amplitude & 1.42     & 1.38     \\
Std deviation  & 0.87     & 0.91     \\
\hline
\end{tabular}
\end{frame}

% -----------------------------------------------------------------------
% Conclusions
% -----------------------------------------------------------------------
\section{Conclusions}

\begin{frame}{Summary}
\begin{itemize}
  \item All conversion features operate as expected \citep{madden1971}
  \item The pipeline handles \alert{edge cases} cleanly
  \item Ready for production use
\end{itemize}
\end{frame}

\clearpage

% -----------------------------------------------------------------------
% AGU endmatter sentinels
% (compiled as normal Beamer frames in slides; extracted for manuscript)
% -----------------------------------------------------------------------
%% AGU_OPENRESEARCH_BEGIN
\begin{frame}{Open Research}
\begin{itemize}
  \item Forecast data archived at \url{https://zenodo.org/record/test}
  \item Analysis code at \url{https://github.com/test/mjo-analysis}
\end{itemize}
\end{frame}
%% AGU_OPENRESEARCH_END

%% AGU_COI_BEGIN
\begin{frame}{Conflict of Interest}
The authors declare no conflicts of interest relevant to this study.
\end{frame}
%% AGU_COI_END

%% AGU_ACKS_BEGIN
\begin{frame}{Acknowledgments}
Supported by NSF grant AGS-0000000. The authors thank the NeuralGCM
development team.
\end{frame}
%% AGU_ACKS_END

\bibliographystyle{plainnat}
\bibliography{refs}

% -----------------------------------------------------------------------
% Supplemental  (AGU: moved after \bibliography; SI checklist counts from here)
% Expected counts: 2 figures, 1 table  → "Figures S1 to S2", "Table S1"
% -----------------------------------------------------------------------
\section{Supplemental}

\begin{frame}{SI: Extended EOF Patterns}
\fig{test_si_figure_1.pdf}
\end{frame}

\begin{frame}{SI: Sensitivity Analysis}
\fig{test_si_figure_2.pdf}
\end{frame}

\begin{frame}{SI: Full Statistics}
\begin{tabular}{lccc}
\hline
Region    & N    & Mean & Std  \\
\hline
Tropics   & 1388 & 1.42 & 0.87 \\
IndOcean  & 1388 & 1.18 & 0.73 \\
WestPac   & 1388 & 1.31 & 0.81 \\
\hline
\end{tabular}
\end{frame}

\end{document}
